% !TeX root = ../papers/nc-p3.tex
\section{第三長方形の全点の無印同型と周期}
\label{sec:third-rectangle-orbits}
\begin{theorem}\label{thm:third-rectangle-orbits}
第三長方形全体で $v=A,h=AD$ と置く。指定始点での完全不変量は
\[
 (X,Y,h)=\bigl(v^2,(1-2/b)/v,h\bigr),\qquad Xh(h-1)(1-XY^2)\ne0
\]
であり、この開集合の各点は一つの窓の同型類を与える。
二位置移動は $(X,Y,h)\mapsto(h^2/X,XY/h,h)$ である。
最小周期は $X=h$ または $(X=-h,Y=0)$ のとき2、他は4である。
一零条件の台89の閉包は $(h-2)^2=h^2XY^2$、その二零条件の台88は $h=2,Y=0$ である。
台89は既約で、その閉包は台88に沿って局所エタール位相で二枚が横断する。
\end{theorem}
\begin{proof}[全係数での四座標枠の復元]
$G$ を $W$ の四次交代収縮の勾配とする。$D_4(W)$ では割らずに使う。
$W_{pq}$ を第 $p,q$ 位置を出力とする切断とし、
$T_{pq}=\operatorname{tr}_q((W_{pq}G_{pq}^t)^2)$ とする。
これは基底変更で共役と共通の指標倍を受ける。四元群作用により対角である。
さらにある位置で対角な共変作用素を挿入して $W$ と $G$ を他の位置で縮約できる。
入力側の作用素が一つの場合、位置 $p$ の第 $i$ 対角成分は
$\sum_{k_p=i}W_kG_k T(k_q)$ である。既に復元した位置では、固有直線に沿う射影子も入力に使える。

位置番号を $0,1,2,3$ とする。まず位置1で $T_{10},T_{12},T_{13}$ と、位置0の $T_{01}$ を挿入した縮約を取る。
次に位置3で、位置1の四射影を挿入した四式と $T_{30},T_{31}$ を使う。
最後に位置0,2で、位置1,3の四射影の組を挿入した16式を取る。
各段階、四文字のどの一対でも同時固有値の全差は単位イデアルを生成する。
これは次の小さい有理係数環計算で確認できる。各切断を総色で四ブロックに分け、
$G_k$ を各位置で $k_j$ を除いた三文字の交代和で作る。
$U=vh(h-1)b(b-1)$ とし、$W,G$ の共通因子を割る際には全て $U$ の因子か非零定数であることを確かめる。
各文字対で差の分子を順に加え、既に得たGröbner基底による剰余とその基底自体から $U$ の因子を除く。
この操作で $\mathbb Q[v,h,b,U^{-1}]$ での生成イデアルは変わらない。
位置 $1,3,0,2$ の順にそれぞれ六対とも基底 $\{1\}$ を得る。
従って全幾何点で同時固有直線が復元でき、任意の無印同型が四位置別の単項行列になる。
\path{computations/verify_third_rectangle_frames.py} は差の生成から24回の単位判定までを行う。
\end{proof}

\begin{proof}[台の比較と商]
始点移動後の台と目標の台が同じ大きさのとき、第一から第三位置までの置換を全て選ぶ。
各語の総色零から第四位置の像が強制され、これが置換で実際に目標の台を与える場合を残す。
非零の薄い語は全て薄い語に移る。薄い係数を1に保つ残余対角作用は厚い最初の八成分に一つの倍率、
後の八成分にその逆倍率を掛けるので、$x_2=1$ で再規格化できる。
零条件なしと二零条件では、元の始点で恒等と
\[
 (v,h,b)\longmapsto(-v,h,b/(b-1))
\]
の二つだけを得る。一零条件では恒等が同じ零位置を、後の変換が反対の零位置を与える。
一・三位置移動では台を移す置換がなく、二位置移動では上の結果に $v\mapsto h/v$ を合成する。
\path{computations/verify_third_rectangle_orbits.py} が台ごとの全置換と係数の再規格化を検算する。

$z=1-2/b$ は $z\ne\pm1$ を動き、上の対合は $(v,z)\mapsto(-v,-z)$ となる。
従って $X=v^2,Y=z/v$ が商の座標を与え、二位置移動の固定条件から周期式を得る。
$h=b$ と $h=b/(b-1)$ の二面は交換され、それぞれ既約で $h=b=2$ で横断する。
二重被覆は $v\ne0$ でエタールであり、商に降ろすと定理の方程式になる。
\end{proof}

\begin{corollary}
$r=v+h/v$ と置く。反対階数は16から
$b(r-2)=h-1$ と $b(r+h+1)=h-1$ のうち成立する式の数を引いたものである。
階数14は $h=-3$ かつ $b(v-3)(v+1)=-4v$ の曲線で尽きる。
台88の全点は三変数族全体における標準AS正則開集合に含まれ、特に台89の一般点も正則である。
\end{corollary}
\begin{proof}
反対切断の色2,3の行列式はともに $-1$ である。色0,1には定数単元の三次小行列式があり、
行列式は順に $-[b(r-2)-h+1]/((b-1)(h-1))$ と $[b(r+h+1)-h+1]/(h-1)$ となる。
後二式の同時消滅では $b\ne0$ から $h=-3$ を得る。
台88では前節の全点検算により五次まで標準であり、必要な二文字の四次中心核は次元3,2を持ち、
その中に長さ13で商を零にする四式がある。軌道を移すときもこの二文字は保たれる。
これらの次元は三変数族の一般値である。各点の近傍で五次までの経路と二つの核は局所自由だから、
この四式を局所的に延長できる。中心商の有限個の次元上界と一般点からの全次数の下界に
定理\ref{thm:parallel-third-rectangle}の証明を適用する。
\end{proof}
